\documentclass[11pt,a4paper]{article}
\usepackage[T1]{fontenc}
\usepackage[utf8]{inputenc}
\usepackage[polish]{babel}
\usepackage{lmodern}
\usepackage{microtype}
\usepackage[a4paper,margin=2.3cm]{geometry}
\usepackage{amsmath,amssymb}
\usepackage{booktabs,longtable}
\usepackage{xcolor}
\usepackage{hyperref}
\usepackage{listings}
\usepackage{enumitem}

\hypersetup{colorlinks=true,linkcolor=blue,urlcolor=blue}
\lstset{basicstyle=\ttfamily\small,breaklines=true,frame=single}

\title{Smog AI HF20: kontrakt czasu prognozy, walidacja opadu i MLOps}
\author{Edward Szczypka}
\date{\today}

\begin{document}
\maketitle

\section{Rozdzielenie czasu serwowania i czasu modelu}

Niech $t_s$ oznacza ostatnią wspólną godzinę danych źródłowych,
$t_c$ chwilę utworzenia prognozy, a $t_a$ najbliższą pełną godzinę ściśle po
$t_c$. Dla leadu serwującego $\ell\in\{1,\ldots,48\}$:
\[
 t_{\ell}=t_a+(\ell-1)\mathrm{h}.
\]
Horyzont modelu wynosi:
\[
 h_m(\ell)=\frac{t_{\ell}-t_s}{1\mathrm{h}}.
\]
Przy maksymalnym wieku źródła równym $12$ godzin otrzymujemy:
\[
 h_{m,\max}=48+12=60.
\]

Pole \texttt{forecast\_horizon} ma znaczenie leadu użytkownika $1\ldots48$,
natomiast \texttt{model\_horizon\_hours} jest cechą modelu i przy opóźnieniu
źródła może należeć do zakresu $5\ldots52$ albo $13\ldots60$.

\section{Warunki odrzucenia predykcji}

Predykcja zostaje odrzucona, jeżeli:
\begin{enumerate}[nosep]
  \item wiek danych przekracza skonfigurowany próg;
  \item aktywny artefakt nie deklaruje wszystkich wymaganych horyzontów;
  \item horyzont modelu przekracza $60$ godzin;
  \item czas docelowy nie jest przyszłą pełną godziną;
  \item siatki czasu poszczególnych parametrów nie są identyczne.
\end{enumerate}

\section{Próbkowanie treningowe}

Dane surowe nie są redukowane. Redukcja dotyczy wyłącznie materializowanej
próby treningowej. Koszyki horyzontów:
\[
[1,6],\ [7,12],\ [13,24],\ [25,48],\ [49,60].
\]
W każdym koszyku wybierana jest ograniczona liczba horyzontów dla danego czasu
bazowego, lecz globalnie wszystkie horyzonty są reprezentowane.

\section{Model hurdle opadu}

Niech $Z=\mathbb{1}\{Y>\tau\}$ oznacza wystąpienie opadu, a $A=Y\mid Z=1$
ilość opadu w okresie mokrym. Model zwraca:
\[
 \widehat p=P(Z=1\mid X),\qquad
 \widehat a=\mathbb{E}[A\mid Z=1,X],\qquad
 \widehat y=\widehat p\widehat a.
\]

Brier Skill Score względem baseline'u $b$:
\[
 \operatorname{BSS}_b=1-\frac{\operatorname{BS}_{model}}
 {\operatorname{BS}_b}.
\]
Brama wymaga dodatniego BSS względem klimatologii i persistence, poprawy MAE,
$\operatorname{AUC}\ge0.60$ oraz kontrolowanego biasu. Niespełnienie dowolnego
warunku nadaje status \texttt{experimental} i blokuje publikację.

\section{MLflow}

MLflow jest opcjonalnym Bridge. Lokalny backend:
\begin{lstlisting}
C:\ProgramData\SmogAI\mlflow\mlflow.db
C:\ProgramData\SmogAI\mlflow\artifacts
http://127.0.0.1:5000
\end{lstlisting}
Rejestrowane są metryki kandydatów, konfiguracja próby, horyzonty,
\texttt{dataset\_id}, SHA-256 i wybór modelu. Domyślny App Spec nie tworzy
serwera MLflow w DigitalOcean, ponieważ oznaczałoby to dodatkową usługę i koszt.

\section{Langfuse}

Tracing jest domyślnie wyłączony. Po jawnej aktywacji każdy prompt może otrzymać
\texttt{trace\_id}, wersję szablonu, metadane modeli i ocenę użytkownika. Bez
Langfuse oceny są zapisywane lokalnie w append-only JSONL.

\section{Polityka publikacji}

Publikacja wymaga jawnej zgody. Do ObjectStore mogą trafić jedynie:
\[
\{\text{model},\text{karta},\text{metryki},\text{aktywny pointer},
\text{porównanie}\}.
\]
Zbiór zabroniony:
\[
\{\text{raw data},\text{SQLite},\text{snapshot},\text{training frames},
\text{source cache}\}.
\]

\section{Architektura DigitalOcean}

\begin{lstlisting}
local Windows pipeline
  -> explicit approved artifacts -> DigitalOcean Spaces
  -> FastAPI read-only service -> Streamlit dashboard
\end{lstlisting}
App Platform nie wykonuje treningu ani backfillu. Opcjonalny MLflow i Langfuse
są niezależnymi decyzjami wdrożeniowymi.

\end{document}
